\section{Limitations and Conclusion}\label{sec:conclusion}
\subsection{Threats to Validity}
\emph{Internal validity:} Random allocation reduces systematic baseline differences, but a small sample does not guarantee balanced spatial ability, attention, or prior navigation experience. The MAP and TBT conditions jointly differ in information format, access time, and whether guidance is available during walking; their individual contributions cannot be identified. Unrecorded route deviations or distractions could further influence scores. \emph{Construct validity:} Five direction choices and ten recognition photographs measure specific aspects of route memory, not complete survey knowledge or long-term retention. A two-minute delay measures short-term performance. \emph{External validity:} One route, a convenience sample of young adults, and a single traversal limit generalization to other ages, environments, devices, and repeated trips. \emph{Statistical conclusion validity:} With ten observations per condition, bootstrap CIs are approximate and potentially sensitive to sampling variability. A positive CI for landmark recognition is evidence within this dataset, not a guarantee of replication.

\subsection{Conclusions and Future Work}
For RQ1, participants reported more frequent use of TBT than prior digital map study, a descriptive finding rather than an experimental effect. For RQ2, MAP achieved higher mean direction recall (78\% vs. 72\%), but its interval includes zero and the revised one-sided test was not significant ($p=0.2894$). For RQ3, MAP achieved higher landmark recognition (85\% vs. 67\%), with a positive bootstrap interval and a significant revised one-sided test ($p=0.00065$). Together, the findings provide preliminary support for the predicted MAP advantage, most clearly for landmark recognition, without establishing superiority across all spatial-memory tasks. Future studies should recruit larger samples, use multiple routes, measure longer-term retention, and manipulate guidance format and timing separately. These refinements would help distinguish mechanisms of learning from the practical effectiveness of the two complete navigation protocols.
